\maketitle

\section{Keller Groupoids: Current Theory
Snapshot}\label{keller-groupoids-current-theory-snapshot}

\textbf{Snapshot:} 2026-09-14\\
\textbf{Purpose:} freeze the current mathematical state before the next
proof phase.

This document is intentionally more conservative than a paper draft.
Every item is separated into one of four provenance/status classes:

\begin{itemize}
\tightlist
\item
  \textbf{published input} --- a theorem imported from the literature;
\item
  \textbf{derived / LLM proved} --- a complete argument has been written
  in the project, but has not yet received independent human
  verification;
\item
  \textbf{human verified} --- explicitly checked by a human;
\item
  \textbf{open} --- conjectural, incomplete, or still under audit.
\end{itemize}

At this snapshot no new Keller-groupoid theorem is marked human verified
or Lean verified.

The central new objective is to turn the Keller-groupoid /
fiber-configuration formalism into constraints on hypothetical plane
Keller maps. The strongest current chain reduces possible failure of
étale-maximality to a remarkably small plane-curve configuration.

\begin{center}\rule{0.5\linewidth}{0.5pt}\end{center}

\section{1. Basic objects}\label{basic-objects}

Let

\[
F:\mathbb A^n_{\mathbb C}\longrightarrow \mathbb A^n_{\mathbb C}
\]

be a polynomial map with nonzero constant Jacobian determinant. Such a
map is called a \textbf{Keller map}. It is étale.

Let \(d\) denote its generic degree.

For \(r\ge1\), define the \(r\)-fold self-fiber product

\[
K_r(F)
=
\underbrace{\mathbb A^n\times_F\cdots\times_F\mathbb A^n}_{r\text{ factors}}
=
\{(x_1,\ldots,x_r):F(x_1)=\cdots=F(x_r)\}.
\]

The spaces \(K_r(F)\) form the Čech nerve of the étale map from the
source to its image.

Define the ordered distinct-point fiber configuration variety

\[
\operatorname{Conf}_k(F)
=
\{(x_1,\ldots,x_k)\in K_k(F):x_i\ne x_j\text{ for }i\ne j\}.
\]

We reserve \(C,C_i,\overline C\) for plane curves.

The relation groupoid

\[
K_2(F)\rightrightarrows K_1(F)=\mathbb A^n
\]

is the \textbf{Keller groupoid}.

\begin{center}\rule{0.5\linewidth}{0.5pt}\end{center}

\section{2. General Keller-groupoid
theorems}\label{general-keller-groupoid-theorems}

\subsection{Theorem KG-G1 --- smooth affine parallelizable
levels}\label{theorem-kg-g1-smooth-affine-parallelizable-levels}

\textbf{Status:} derived / LLM proved.

Every \(K_r(F)\) and every open-and-closed configuration stratum
\(\operatorname{Conf}_k(F)\) is a smooth affine \(n\)-fold. Every
projection to one source coordinate is étale, hence

\[
T_{K_r(F)}\simeq\mathcal O^{\oplus n},
\qquad
T_{\operatorname{Conf}_k(F)}\simeq\mathcal O^{\oplus n}.
\]

Thus every fiber configuration variety is algebraically parallelizable.

\begin{center}\rule{0.5\linewidth}{0.5pt}\end{center}

\subsection{Theorem KG-G2 --- partition
decomposition}\label{theorem-kg-g2-partition-decomposition}

\textbf{Status:} derived / LLM proved.

For every \(r\), equality of coordinates defines an open-and-closed
decomposition indexed by set partitions \(\pi\) of \(\{1,\ldots,r\}\):

\[
K_r(F)
\simeq
\coprod_{\pi\in\Pi(r)}
\operatorname{Conf}_{|\pi|}(F).
\]

Equivalently,

\[
K_r(F)
\simeq
\coprod_{k\ge1}
S(r,k)\,\operatorname{Conf}_k(F),
\]

where \(S(r,k)\) is a Stirling number of the second kind.

\begin{center}\rule{0.5\linewidth}{0.5pt}\end{center}

\subsection{Theorem KG-G3 --- finite configuration
palette}\label{theorem-kg-g3-finite-configuration-palette}

\textbf{Status:} derived / LLM proved.

No geometric fiber has more than \(d\) points. Therefore

\[
\operatorname{Conf}_k(F)=\varnothing\qquad(k>d).
\]

Consequently the infinite Čech nerve has only the finite geometric
palette

\[
\operatorname{Conf}_1(F),\ldots,\operatorname{Conf}_d(F).
\]

Every connected-component isomorphism type occurring anywhere in the
nerve occurs by level \(d\).

\begin{center}\rule{0.5\linewidth}{0.5pt}\end{center}

\subsection{Theorem KG-G4 --- fiber-cardinality
filtration}\label{theorem-kg-g4-fiber-cardinality-filtration}

\textbf{Status:} derived / LLM proved.

Let

\[
U_{\ge k}
=
\{y:\#F^{-1}(y)\ge k\}.
\]

The common-image map

\[
q_k:\operatorname{Conf}_k(F)\to\mathbb A^n
\]

is étale and has image \(U_{\ge k}\).

Thus the configuration tower geometrizes the exact fiber-cardinality
filtration.

Over an exact \(m\)-point fiber the map \(q_k\) has fiber cardinality

\[
(m)_k=m(m-1)\cdots(m-k+1).
\]

\begin{center}\rule{0.5\linewidth}{0.5pt}\end{center}

\subsection{Theorem KG-G5 --- top forgetful
map}\label{theorem-kg-g5-top-forgetful-map}

\textbf{Status:} derived / LLM proved.

The forgetful morphism

\[
\rho_d:
\operatorname{Conf}_d(F)\longrightarrow
\operatorname{Conf}_{d-1}(F)
\]

is an open immersion.

It is an isomorphism exactly when no fiber of \(F\) has cardinality
\(d-1\).

\begin{center}\rule{0.5\linewidth}{0.5pt}\end{center}

\subsection{Theorem KG-G6 --- monodromy orbit
theorem}\label{theorem-kg-g6-monodromy-orbit-theorem}

\textbf{Status:} derived from standard finite-cover theory / LLM proved.

Let

\[
V=U_{\ge d}
\]

be the proper finite-étale locus and let

\[
G\le S_d
\]

be the monodromy group.

Then over \(V\),

\[
\operatorname{Conf}_k(F)|_V
\]

is the finite étale cover associated to the \(G\)-set of ordered
injective \(k\)-tuples of sheets.

Connected components correspond to \(G\)-orbits on ordered distinct
\(k\)-tuples. In particular, \(\operatorname{Conf}_k(F)\) is generically
connected exactly when the monodromy action is \(k\)-transitive.

\begin{center}\rule{0.5\linewidth}{0.5pt}\end{center}

\subsection{Corollary KG-G7 --- top configuration and Galois
closure}\label{corollary-kg-g7-top-configuration-and-galois-closure}

\textbf{Status:} derived / LLM proved.

Over \(V\), \(\operatorname{Conf}_d(F)\) is the full frame torsor of the
cover.

It has

\[
\frac{d!}{|G|}
\]

connected components, each realizing the Galois closure with deck group
\(G\).

\begin{center}\rule{0.5\linewidth}{0.5pt}\end{center}

\subsection{Theorem KG-G8 --- additive/Stirling
calculus}\label{theorem-kg-g8-additivestirling-calculus}

\textbf{Status:} derived / LLM proved.

For every additive invariant \(I\),

\[
I(K_r(F))
=
\sum_{k=1}^d S(r,k) I(\operatorname{Conf}_k(F)).
\]

Stirling inversion recovers the configuration invariants from the
nerve-level invariants.

The scalar sequence \(I(K_r(F))\) satisfies a linear recurrence whose
characteristic polynomial divides

\[
\prod_{j=1}^d(T-j).
\]

\begin{center}\rule{0.5\linewidth}{0.5pt}\end{center}

\section{3. Euler and fiber-deficit
identities}\label{euler-and-fiber-deficit-identities}

Let

\[
U_m=\{y:\#F^{-1}(y)=m\}.
\]

\subsection{Theorem KG-E1 --- factorial-moment
identity}\label{theorem-kg-e1-factorial-moment-identity}

\textbf{Status:} derived / LLM proved.

For every \(k\),

\[
\boxed{
\chi(\operatorname{Conf}_k(F))
=
\sum_{m=k}^d (m)_k\,\chi(U_m).
}
\]

Thus the configuration Euler characteristics are factorial moments of
the fiber-cardinality stratification.

The system is triangular and recovers \(\chi(U_m)\) from the
configuration Euler characteristics.

\begin{center}\rule{0.5\linewidth}{0.5pt}\end{center}

\subsection{Theorem KG-E2 --- global Euler-deficit
identity}\label{theorem-kg-e2-global-euler-deficit-identity}

\textbf{Status:} derived / LLM proved; high-priority human audit.

For a plane Keller map of generic degree \(d\),

\[
\boxed{
d-1
=
\sum_{m=0}^{d-1}(d-m)\,\chi(U_m).
}
\]

\subsubsection{Derivation}\label{derivation}

Euler integration of fiber cardinality gives

\[
\chi(\mathbb A^2_{\rm source})
=
\sum_{m=0}^d m\,\chi(U_m)
=
1.
\]

Additivity on the target gives

\[
\sum_{m=0}^d\chi(U_m)=1.
\]

Subtracting the first identity from \(d\) times the second gives the
result.

This is one of the main bridges from the Keller groupoid to plane-curve
topology.

\begin{center}\rule{0.5\linewidth}{0.5pt}\end{center}

\section{4. Plane configuration
criteria}\label{plane-configuration-criteria}

Fix now a plane Keller map

\[
F:\mathbb A^2\to\mathbb A^2.
\]

Let \(\overline X\to\mathbb A^2\) be the finite normalization supplied
by Zariski Main.

For an irreducible component \(C\) of the nonproperness curve, let

\begin{itemize}
\tightlist
\item
  \(\sigma_C\in S_d\) be generic inertia;
\item
  \(m_C\) be the generic affine fiber cardinality along \(C\);
\item
  \(t_C\) be the number of generically unramified normalization sheets
  omitted from the affine source.
\end{itemize}

\subsection{Theorem KG-P1 --- deficit
decomposition}\label{theorem-kg-p1-deficit-decomposition}

\textbf{Status:} derived / LLM proved.

\[
\boxed{
d-m_C
=
\operatorname{supp}(\sigma_C)+t_C.
}
\]

Thus fiber deficit decomposes into ramification loss plus unramified
deletion.

The configuration tower recovers \(m_C\), so once inertia is known,
\(t_C\) is derived rather than free data.

\begin{center}\rule{0.5\linewidth}{0.5pt}\end{center}

\subsection{Corollary KG-P2 --- deficit one detects
non-maximality}\label{corollary-kg-p2-deficit-one-detects-non-maximality}

\textbf{Status:} derived / LLM proved.

A nonidentity permutation cannot move exactly one sheet. Hence a generic
\((d-1)\)-point fiber along a curve can only arise from one deleted
unramified sheet:

\[
d-m_C=1
\Longrightarrow
\operatorname{supp}(\sigma_C)=0,\quad t_C=1.
\]

\begin{center}\rule{0.5\linewidth}{0.5pt}\end{center}

\subsection{Theorem KG-P3 --- top stabilization under
étale-maximality}\label{theorem-kg-p3-top-stabilization-under-uxe9tale-maximality}

\textbf{Status:} derived / LLM proved.

If the Keller map is étale-maximal, then it has no \((d-1)\)-point fiber
and

\[
\boxed{
\operatorname{Conf}_d(F)
\simeq
\operatorname{Conf}_{d-1}(F).
}
\]

The degree-three marked-cubic equality
\(\operatorname{Conf}_3\simeq\operatorname{Conf}_2\) is an instance of
this general phenomenon.

\begin{center}\rule{0.5\linewidth}{0.5pt}\end{center}

\subsection{Theorem KG-P4 --- configuration unit-rank
bound}\label{theorem-kg-p4-configuration-unit-rank-bound}

\textbf{Status:} derived / LLM proved.

Let \(Z\) be a connected component of \(\operatorname{Conf}_k(F)\)
dominating its image. If the codimension-one complement of that image
has \(c_k\) irreducible components, then

\[
\operatorname{rank}
\left(\mathcal O(Z)^*/\mathbb C^*\right)
\ge c_k.
\]

For a top-configuration component this gives at least as many
independent units as there are irreducible components of the
nonproperness curve.

\begin{center}\rule{0.5\linewidth}{0.5pt}\end{center}

\subsection{Theorem KG-P5 --- secant
projector}\label{theorem-kg-p5-secant-projector}

\textbf{Status:} derived independently here, but priority overlaps
external 2026 work.

For \(F=(P,Q)\), choose a polynomial divided-difference matrix
\(D_F(x,y)\) satisfying

\[
F(x)-F(y)=D_F(x,y)(x-y)
\]

and write

\[
c=\det JF,\qquad \delta_F=\det D_F.
\]

On \(K_2(F)\),

\[
e_\Delta=\delta_F/c
\]

is the idempotent of the diagonal component. Hence

\[
\delta_F(\delta_F-c)
\in
(P(x)-P(y),Q(x)-Q(y)),
\]

and the off-diagonal configuration factor is cut out by

\[
(P(x)-P(y),Q(x)-Q(y),\delta_F).
\]

Therefore

\[
F\text{ invertible}
\Longleftrightarrow
(P(x)-P(y),Q(x)-Q(y),\delta_F)=(1).
\]

\textbf{Provenance rule:} if this criterion is load-bearing, cite
\emph{Collision Ideals and Off-Diagonal Sheets} at its pinned public
commit and contact Roy van Rijn regarding his earlier independent
off-diagonal/saturation constructions.

\begin{center}\rule{0.5\linewidth}{0.5pt}\end{center}

\section{5. Published plane-curve
inputs}\label{published-plane-curve-inputs}

The following are literature results, not claims of this project.

\subsection{Published input L1 --- Jelonek parametric
components}\label{published-input-l1-jelonek-parametric-components}

Every irreducible component of the nonproperness set of a generically
finite polynomial map \(\mathbb C^2\to\mathbb C^2\) is polynomially
parametric.

\begin{center}\rule{0.5\linewidth}{0.5pt}\end{center}

\subsection{Published input L2 --- Chau one point at
infinity}\label{published-input-l2-chau-one-point-at-infinity}

For a nonsingular polynomial map \(\mathbb C^2\to\mathbb C^2\), a
nonempty nonproperness curve has one projective point at infinity.

Nguyen Van Chau, \emph{Note on the Jacobian condition and the non-proper
value set}, Ann. Polon. Math. 84 (2004), 203--210.

\begin{center}\rule{0.5\linewidth}{0.5pt}\end{center}

\subsection{Derived consequence L2a --- one place per irreducible
component}\label{derived-consequence-l2a-one-place-per-irreducible-component}

\textbf{Status:} derived / LLM proved.

Polynomial parametrization lifts to the normalization. The completed
normalization is \(\mathbb P^1\), and only the source point \(\infty\)
can map to the projective boundary.

Hence each irreducible component has normalization \(\mathbb A^1\) and
exactly one place at infinity.

\begin{center}\rule{0.5\linewidth}{0.5pt}\end{center}

\subsection{Published input L3 --- Jelonek affine-extension boundary
theorem}\label{published-input-l3-jelonek-affine-extension-boundary-theorem}

For a normal affine extension of a smooth acyclic surface, horizontal
boundary components are homeomorphic to \(\mathbb C\) and distinct
horizontal components are disjoint.

Z. Jelonek, \emph{Testing sets for properness of polynomial mappings},
Math. Ann. 315 (1999), 1--35, Theorem 4.6.

\begin{center}\rule{0.5\linewidth}{0.5pt}\end{center}

\subsection{Published input L4 --- connected nonproperness Euler
positivity}\label{published-input-l4-connected-nonproperness-euler-positivity}

If \(S_F\) is connected for a nonsingular polynomial self-map of affine
space, then

\[
\chi(S_F)>0.
\]

Z. Jelonek, \emph{A note on the Jacobian Conjecture}, Colloq. Math. 170
(2022), 85--90.

\begin{center}\rule{0.5\linewidth}{0.5pt}\end{center}

\subsection{Published input L5 --- Assi rational one-place pencil
theorem}\label{published-input-l5-assi-rational-one-place-pencil-theorem}

For a rational polynomial with one place at infinity, either the
polynomial is equivalent to a coordinate and every member of the pencil
is rational, or the pencil has at most two rational members.

A. Assi, \emph{Rational curves with one place at infinity}, 2012.

\begin{center}\rule{0.5\linewidth}{0.5pt}\end{center}

\subsection{Published input L6 --- Assi disjoint-parametric-curve
trichotomy}\label{published-input-l6-assi-disjoint-parametric-curve-trichotomy}

For distinct monic polynomially parametrized plane curves \(f=0,g=0\),
either

\begin{enumerate}
\def\labelenumi{\arabic{enumi}.}
\tightlist
\item
  they are translates in a coordinate pencil;
\item
  they are the two exceptional rational members of a noncoordinate
  one-place pencil; or
\item
  they intersect in the affine plane.
\end{enumerate}

In the exceptional two-fiber case, the two fibers have equal numbers of
singular points; every singular point has two branches and Milnor number
one, hence is an ordinary transverse node.

This is Assi Proposition 3.5 together with Proposition 3.3 and Lemma
3.2.

\begin{center}\rule{0.5\linewidth}{0.5pt}\end{center}

\subsection{Published input L7 --- Braun--Dias--Venato-Santos
intersection
theorem}\label{published-input-l7-braundiasvenato-santos-intersection-theorem}

For a counterexample to the Jacobian conjecture, the nonproperness set
meets every hypersurface

\[
Z=h^{-1}(0)
\]

biregular to \(\mathbb C^{n-1}\) when \(h\) is a polynomial submersion.

In dimension two, a finite union of fibers of a coordinate polynomial is
therefore impossible: choose another parallel fiber disjoint from the
finite union.

\begin{center}\rule{0.5\linewidth}{0.5pt}\end{center}

\section{6. Étale-maximality
reduction}\label{uxe9tale-maximality-reduction}

Let

\[
\mathbb A^2=U\hookrightarrow\overline X
\xrightarrow{\pi}\mathbb A^2
\]

be the finite Zariski-Main normalization.

The Keller map is \textbf{étale-maximal} when \(U\) is the entire étale
locus of \(\pi\).

\begin{center}\rule{0.5\linewidth}{0.5pt}\end{center}

\subsection{Theorem EM-1 --- boundary
purity}\label{theorem-em-1-boundary-purity}

\textbf{Status:} derived / LLM proved.

The boundary

\[
D=\overline X\setminus U
\]

has no isolated points; it is pure divisorial.

\begin{center}\rule{0.5\linewidth}{0.5pt}\end{center}

\subsection{Theorem EM-2 --- generically unramified boundary is
everywhere
étale}\label{theorem-em-2-generically-unramified-boundary-is-everywhere-uxe9tale}

\textbf{Status:} derived / LLM proved.

Let \(E\subset D\) be a boundary prime divisor generically unramified
for \(\pi\). Then \(\pi\) is étale at every point of \(E\).

The proof uses purity of the branch locus plus Jelonek's disjointness of
horizontal boundary components: any point of ramification on \(E\) would
force a different ramification boundary curve to intersect \(E\).

\begin{center}\rule{0.5\linewidth}{0.5pt}\end{center}

\subsection{Theorem EM-3 --- deleted unramified sheet is a normalization
copy}\label{theorem-em-3-deleted-unramified-sheet-is-a-normalization-copy}

\textbf{Status:} derived / LLM proved.

Such an \(E\) is isomorphic to \(\mathbb A^1\). If \(C=\pi(E)\), then

\[
E\simeq\widetilde C\simeq\mathbb A^1
\]

and \(E\to C\) is the normalization morphism.

\begin{center}\rule{0.5\linewidth}{0.5pt}\end{center}

\subsection{Theorem EM-4 --- unibranch singularities cannot support
deletion}\label{theorem-em-4-unibranch-singularities-cannot-support-deletion}

\textbf{Status:} derived / LLM proved.

If \(E\) is an omitted unramified normalization sheet, every singularity
of \(C=\pi(E)\) met by that sheet must be multibranch.

A unibranch singularity would force the normalization map to fail to be
immersive, contradicting étaleness of the ambient surface map along
\(E\).

\begin{center}\rule{0.5\linewidth}{0.5pt}\end{center}

\subsection{Corollary EM-5 --- unibranch
criterion}\label{corollary-em-5-unibranch-criterion}

\textbf{Status:} derived / LLM proved.

If every irreducible component of \(S_F\) is unibranch at every finite
singular point, then \(F\) is étale-maximal.

\begin{center}\rule{0.5\linewidth}{0.5pt}\end{center}

\subsection{Theorem EM-6 --- conductor-point fiber
bound}\label{theorem-em-6-conductor-point-fiber-bound}

\textbf{Status:} derived / LLM proved.

Suppose an omitted unramified normalization sheet lies over a point
\(p\in C\) having \(r_p\) branches. Then

\[
\#F^{-1}(p)\le d-r_p.
\]

If \(t\) omitted unramified normalization sheets lie over the same
component, then

\[
\#F^{-1}(p)\le d-t\,r_p.
\]

This converts conductor self-identification into a configuration-visible
fiber deficit.

\begin{center}\rule{0.5\linewidth}{0.5pt}\end{center}

\section{7. Connectedness obstruction to
non-maximality}\label{connectedness-obstruction-to-non-maximality}

\subsection{Theorem EM-7 --- connected nonproperness implies
étale-maximality}\label{theorem-em-7-connected-nonproperness-implies-uxe9tale-maximality}

\textbf{Status:} derived / LLM proved; high-priority human audit.

If \(S_F\) is connected, then \(F\) is étale-maximal.

\subsubsection{Argument}\label{argument}

Assume non-maximality. By EM-3--EM-4, some irreducible component \(C\)
has normalization \(\mathbb A^1\) and at least one multibranch
self-identification.

Let \(C_1,\ldots,C_c\) be the irreducible components of \(S_F\). Their
normalizations are copies of \(\mathbb A^1\), so

\[
\chi(S_F)
=
c-\sum_{p}(r_p-1),
\]

where \(r_p\) is the total number of normalization branches above \(p\).

Connectedness already requires at least \(c-1\) identifications. A
self-identification on a single irreducible component contributes at
least one additional identification. Therefore

\[
\chi(S_F)\le0.
\]

But Jelonek's published theorem gives

\[
\chi(S_F)>0
\]

when \(S_F\) is connected. Contradiction.

Hence:

\[
\boxed{
F\text{ non-étale-maximal}
\Longrightarrow
S_F\text{ disconnected}.
}
\]

This is currently one of the main new structural outputs.

\begin{center}\rule{0.5\linewidth}{0.5pt}\end{center}

\section{8. Disconnected nonproperness and the Assi
reduction}\label{disconnected-nonproperness-and-the-assi-reduction}

The following chain has been re-audited against Assi's actual
statements.

\subsection{Theorem EM-8 --- disconnected components lie in a common
rational
pencil}\label{theorem-em-8-disconnected-components-lie-in-a-common-rational-pencil}

\textbf{Status:} derived / LLM proved; high-priority human audit.

Assume \(S_F\) is disconnected.

Choose irreducible components \(C\) and \(C'\) in distinct connected
components. They are disjoint in \(\mathbb A^2\), polynomially
parametric, rational, and one-place at infinity.

After the standard monic normalization at infinity, Assi Proposition 3.5
implies that their defining polynomials differ by a constant:

\[
C=V(h-a),\qquad
C'=V(h-b).
\]

Repeating the argument with any other irreducible component and one
fixed component in a different connected component puts every
irreducible component into the same pencil.

Moreover, irreducible components lying in the same connected component
would then be distinct fibers of the same polynomial and hence disjoint.
Therefore each connected component is irreducible.

Thus

\[
\boxed{
S_F=\bigsqcup_{i=1}^q V(h-a_i)
}
\]

for distinct values \(a_i\), with each selected fiber irreducible,
rational, and one-place at infinity.

\begin{center}\rule{0.5\linewidth}{0.5pt}\end{center}

\subsection{Theorem EM-9 --- exactly two selected fibers in the
noncoordinate
case}\label{theorem-em-9-exactly-two-selected-fibers-in-the-noncoordinate-case}

\textbf{Status:} derived / LLM proved; high-priority human audit.

If \(q\ge3\), Assi's rational-pencil theorem forces \(h\) to be
equivalent to a coordinate.

But a finite union of coordinate fibers misses another coordinate fiber.
The Braun--Dias--Venato-Santos intersection theorem says the
nonproperness set of a Keller counterexample must meet every such
\(\mathbb A^1\)-fiber.

Therefore the coordinate alternative is impossible and

\[
\boxed{q=2.}
\]

So every non-étale-maximal planar Keller map would have

\[
S_F=C_0\sqcup C_1
\]

with \(C_0,C_1\) the two exceptional rational members of a noncoordinate
one-place pencil.

\begin{center}\rule{0.5\linewidth}{0.5pt}\end{center}

\subsection{Theorem EM-10 --- two-nodal-fiber
theorem}\label{theorem-em-10-two-nodal-fiber-theorem}

\textbf{Status:} derived / LLM proved; \textbf{highest-priority human
audit}.

If a plane Keller map is not étale-maximal, then its nonproperness curve
has exactly two connected/irreducible components

\[
\boxed{
S_F=C_0\sqcup C_1,
}
\]

and there exists a polynomial \(h\) and distinct \(a_0,a_1\in\mathbb C\)
such that

\[
C_i=V(h-a_i).
\]

Each \(C_i\)

\begin{itemize}
\tightlist
\item
  is rational;
\item
  has one place at infinity;
\item
  has exactly the same positive number \(s\) of affine singular points;
\item
  every affine singularity is an ordinary transverse node.
\end{itemize}

Equivalently, non-maximality forces the plane Keller map into the
exceptional two-rational-fiber case of Assi's classification.

\subsubsection{Provenance}\label{provenance}

The \emph{classification of the two fibers} is Assi's published theorem.
The \emph{implication from Keller non-maximality to that classification
problem} is the new derived chain of this project.

This distinction must be preserved in any paper.

\begin{center}\rule{0.5\linewidth}{0.5pt}\end{center}

\section{9. Consequences of the two-nodal
reduction}\label{consequences-of-the-two-nodal-reduction}

If \(C_i\) has \(s\) nodes and normalization \(\mathbb A^1\), then

\[
\chi(C_i)=1-s,
\qquad
\chi(S_F)=2-2s\le0.
\]

Every omitted unramified sheet over one of the \(C_i\) is the
normalization copy \(\mathbb A^1\to C_i\), and at each node it
contributes two missing points of the finite normalization fiber.

Thus node strata create fiber deficits visible in the configuration
filtration.

This is the next place to combine

\[
d-1=\sum_{m<d}(d-m)\chi(U_m)
\]

with the equal-node-count theorem.

\begin{center}\rule{0.5\linewidth}{0.5pt}\end{center}

\section{10. What is still open}\label{what-is-still-open}

\subsection{Open O1 --- full planar
étale-maximality}\label{open-o1-full-planar-uxe9tale-maximality}

Prove that the two-nodal-fiber configuration of EM-10 cannot occur for a
Keller map.

This would prove that every plane Keller map is étale-maximal.

\begin{center}\rule{0.5\linewidth}{0.5pt}\end{center}

\subsection{Open O2 --- Euler-deficit
contradiction}\label{open-o2-euler-deficit-contradiction}

Compute the exact contribution of

\begin{itemize}
\tightlist
\item
  the generic portions of \(C_0,C_1\);
\item
  their \(2s\) nodal points;
\item
  ramified versus deleted normalization sheets
\end{itemize}

to the global Euler-deficit identity.

Try to show that no nonnegative integer sheet/inertia data satisfy the
resulting equation.

\begin{center}\rule{0.5\linewidth}{0.5pt}\end{center}

\subsection{Open O3 --- delta-sequence / node
budget}\label{open-o3-delta-sequence-node-budget}

Encode each \(C_i\) by its one-place-at-infinity delta-sequence. Use the
genus formula and equal node counts to restrict or enumerate possible
sequences.

\begin{center}\rule{0.5\linewidth}{0.5pt}\end{center}

\subsection{Open O4 --- boundary/Chern
contradiction}\label{open-o4-boundarychern-contradiction}

Lift the two-fiber decorated pencil through the finite normalization,
resolve it, and impose the affine-plane boundary conditions

\[
\det Q=(-1)^{r-1},
\qquad
\operatorname{sign}Q=(1,r-1),
\qquad
K^2=10-r.
\]

\begin{center}\rule{0.5\linewidth}{0.5pt}\end{center}

\subsection{Open O5 --- Lean proof
development}\label{open-o5-lean-proof-development}

Replace the statement stubs in \texttt{lean/} by proofs, allowing
explicitly identified published theorems to remain axioms.

\begin{center}\rule{0.5\linewidth}{0.5pt}\end{center}

\section{11. Explicit degree-three
example}\label{explicit-degree-three-example}

For the marked-cubic Keller map \(F:\mathbb A^3\to\mathbb A^3\):

\[
\operatorname{Conf}_1(F)=\mathbb A^3,
\qquad
\operatorname{Conf}_2(F)=Y,
\qquad
\operatorname{Conf}_3(F)\simeq Y,
\qquad
\operatorname{Conf}_k(F)=\varnothing\quad(k\ge4).
\]

Hence

\[
K_r(F)
\simeq
\mathbb A^3
\sqcup
\frac{3^r-3}{6}\,Y.
\]

This was the motivating example for the Keller-groupoid formalism, but
the planar theory above is now logically independent of the explicit
cubic formulas.

\begin{center}\rule{0.5\linewidth}{0.5pt}\end{center}

\section{12. Citation/provenance
reminders}\label{citationprovenance-reminders}

\begin{enumerate}
\def\labelenumi{\arabic{enumi}.}
\tightlist
\item
  The symmetric-product marked-root interpretation of the explicit
  degree-three example should be attributed according to
  \texttt{CITATION\_HYGIENE.md}.
\item
  If the secant-projector criterion becomes load-bearing, cite the
  pinned public source for \emph{Collision Ideals and Off-Diagonal
  Sheets} and contact Roy van Rijn about the relation to his earlier
  independent off-diagonal/saturation work; see
  \texttt{ROY\_VAN\_RIJN\_PROVENANCE.md}.
\item
  EM-10 must never be described as ``Assi's theorem.'' Assi proves the
  curve-pencil classification. The Keller-to-Assi reduction is our
  derived contribution.
\end{enumerate}

\begin{center}\rule{0.5\linewidth}{0.5pt}\end{center}

\section{13. Perfect monodromy and degree six (post-v11
additions)}\label{perfect-monodromy-and-degree-six-post-v11-additions}

\subsection{Theorem PM-1 --- perfect monodromy in the non-maximal
residual
case}\label{theorem-pm-1-perfect-monodromy-in-the-non-maximal-residual-case}

\textbf{Status:} llm proved / highest-priority human audit.

Assuming the preceding two-nodal-fiber reduction, the sheet-monodromy
group of a non-étale-maximal planar Keller map is perfect: \[
G_{\rm mon}^{\rm ab}=1.
\] The proof uses the nodal vanishing-cycle basis and conjugacy of the
two branch-inertia permutations at each node. See
\texttt{PERFECT\_MONODROMY\_DEGREE6.md} for the load-bearing audit
points.

\subsection{Corollary PM-2 --- S6 is excluded
conditionally}\label{corollary-pm-2-s6-is-excluded-conditionally}

\textbf{Status:} llm proved, conditional on PM-1.

Since \(S_6\) has the nontrivial sign quotient, it is not perfect. Thus
PM-1 excludes \(S_6\) from the residual non-maximal degree-six case.

\subsection{External research input D6-EXT --- current six-sheet
frontier}\label{external-research-input-d6-ext-current-six-sheet-frontier}

\textbf{Status:} external public research result, not a published axiom.

The current independent exact degree-six analysis in
\texttt{alok/jacobian-two} proves that a hypothetical degree-six plane
Keller counterexample has monodromy \(A_6\) or \(S_6\). Combining this
with PM-1 leaves \(A_6\).

\subsection{Theorem A6-1 --- forced A6 boundary
packet}\label{theorem-a6-1-forced-a6-boundary-packet}

\textbf{Status:} llm proved / high-priority audit.

In the residual degree-six \(A_6\) case, the boundary packet is forced
to be \[
(3,1),(1,1),(1,1),
\] with zero excess: one single-3-cycle ramified row and two unramified
deleted rows.

\subsection{Theorem A6-2 --- forced node
inertia}\label{theorem-a6-2-forced-node-inertia}

\textbf{Status:} llm proved / high-priority audit.

At every node of the ramified Assi fiber, the two local branch inertia
permutations are disjoint 3-cycles, hence the local inertia group is
\(C_3\times C_3\).

\subsection{Open A6-O1 --- eliminate the constrained A6
quotient}\label{open-a6-o1-eliminate-the-constrained-a6-quotient}

Rule out a transitive surjection \[
\pi_1(\mathbb A^2\setminus C)\twoheadrightarrow A_6
\] for an Assi exceptional rational nodal one-place curve, with
geometric meridians mapping to single 3-cycles and node branch pairs
mapping to disjoint 3-cycles.

